\chapter{Tilt-Table Test}

\section*{Overview}
A tilt-table test evaluates how the heart rate and blood pressure respond when the body changes from lying flat to an upright position. It is used mainly to investigate fainting or unexplained lightheadedness. The exact approach, setting, and preparation depend on the indication, anatomy, and the patient's health.

\section*{Why It Is Done}
The test may help evaluate suspected vasovagal syncope, orthostatic intolerance, or other disorders of autonomic blood-pressure control when the cause is not clear from the history and routine examination.

\section*{How It Is Performed}
The patient lies on a motorized table secured with straps. Heart rhythm, blood pressure, and symptoms are monitored while the table is slowly tilted upright for a set period. In some protocols, a medicine is used to increase the likelihood of reproducing symptoms.

\section*{Preparation and Risks}
Fasting and temporary changes to selected medicines may be recommended. Follow the testing center's instructions. Dizziness, nausea, sweating, fainting, palpitations, or a temporary blood-pressure change may occur; monitoring and staff support are provided throughout.

\section*{Outcome and Follow-up}
The clinician compares symptoms with changes in blood pressure and heart rate. A positive pattern may support a diagnosis such as vasovagal syncope or orthostatic hypotension, but results must be interpreted with the full clinical history.

\section*{Contraindications and Precautions}
The team reviews unstable heart disease, severe aortic obstruction, uncontrolled arrhythmia, severe anemia or dehydration, pregnancy, and the ability to tolerate upright positioning. The test may be postponed when fainting could cause immediate danger or urgent evaluation is needed first.

\section*{When to Seek Medical Care}
Follow the procedure team's specific instructions. Seek urgent medical assessment for severe or worsening pain, uncontrolled bleeding, fever or signs of infection, trouble breathing, chest pain, fainting, new weakness or confusion, inability to urinate, or any rapidly worsening symptom. The appropriate warning signs depend on the procedure and the patient's medical condition.

\section*{References}
\begin{enumerate}
  \item MedlinePlus. Tilt-table test. U.S. National Library of Medicine; 2025.
  \item Current Medical Diagnosis and Treatment. McGraw-Hill; 2025.
\end{enumerate}
\clearpage
